\section{事業継続：原因より重要事業への影響を調べる}
\label{theory:safety:bcp}
BCPは事業継続計画、BCMはその計画を支える経営・管理の活動である。重要な事業を継続し、又は目標とする期間・水準で復旧させるために、事業への影響、必要な資源、依存関係、代替手段を整理する。自然災害に限らず、感染症、設備事故、サイバー攻撃、取引先の中断等も事業へ影響する\parencite{eth54_bcp2013,eth54_bcp2021}。

復旧が容易な仕事を先に再開すればよいとは限らない。供給停止が顧客、収入、雇用、法令・契約、地域機能へ与える影響と許容時間を評価し、重要業務を選ぶ。バックアップは本体と同じ原因で失われないよう、場所、電源、通信、供給者等の共通の依存を点検する。

\paragraph{計画から実行・改善へ}
経営者が責任を持ち、訓練、点検、資源の確保、事業環境の変化への見直しを進める。計画書を作ったことだけでは対応能力を保証しない。重要業務や依存関係の可視化は平時の改善にもつながる。地域と取引先との協力を含め、組織単独で復旧できると仮定しない。

非常時だから法令や契約がすべて無効になるわけではない。緊急措置には固有の条件があり、権限や責任を確認する。「全て遵守しなくてよい」という包括的免除とは区別する。

\paragraph{普及率は調査年と分母を読む}
平成29年度調査では策定済みの割合が大企業64.0\%、中堅企業31.8\%とされる\parencite{eth54_bcp_survey2018}。これはすべての会社を同じ分母で測った数字ではない。策定中も含む率、企業規模、回答数、調査時点を分ける。令和元年度問題を後に公表された調査で説明しない。

\section{予防的取組：不確実性の中で行動を選ぶ}
\label{theory:safety:precaution}
知識が不完全なことと危険がないことは別である。深刻又は不可逆的な影響が懸念されるときは、科学的確実性がないことだけを理由に合理的な低減策を延期しない。利益、費用、代替案、行動しない場合の損失を比較し、対策の程度と見直しの条件を定める\parencite{eth54_rio,eth54_eea}。

平成25年度問題の十二項目は、無知と不確実性、監視、科学の盲点、学際的学習、実世界条件、利益とリスク、代替策、地域の知識、異なる価値、規制の独立性、制度の障壁、分析だけによる停滞を点検する。これらを考慮することは、科学を捨てることやあらゆる技術を禁止することではない。原典前文の宣言の「批准」という表現も、条約の法的手続と同一にしない。

\section{対話の視点と検証可能性}
\label{theory:safety:riskdialogue}
管理する側の集団的な確率評価と、危険を受ける当事者の判断は異なる。自分で選択できるか、利益と危険の配分は公平か、組織を信頼できるか等もリスク認知に関わる。ハザードとアウトレージを合わせる表現は認知モデルであり、単位の異なる物理量を足す計算式ではない\parencite{eth54_riskcomm}。

専門家の知識の伝達だけで対話が完成するわけではない。中立的な媒介者の支援、権限と情報の非対称性への配慮、異なる根拠の提示、検証可能性が必要になる。分かりやすい説明を作るときにも、情報源や不確かさを隠してはならない。
